\documentclass[11pt]{article}
\usepackage[margin=1in]{geometry}
\usepackage{amsmath,amssymb,amsthm,booktabs}
\usepackage[colorlinks=true,linkcolor=blue,citecolor=blue,urlcolor=blue]{hyperref}
\newtheorem{theorem}{Theorem}
\newcommand{\VR}{\operatorname{VR}_{<}}
\newcommand{\conv}{\operatorname{conv}}
\title{Six points with a contractible Rips shadow and nontrivial fundamental group}
\author{Kaiwen Jin\\\small\texttt{yui.ui.siki@gmail.com}}
\date{13 October 2026}

\begin{document}
\maketitle

\begin{abstract}
The affine projection from a Euclidean Vietoris--Rips complex to its shadow
need not be injective on fundamental groups. We give six rational points in
three-dimensional space whose Rips complex retracts onto a four-cycle, while
its shadow is the union of two convex quadrilaterals containing the origin
and is therefore contractible. Scaling and isometric inclusion give the same
failure in every ambient dimension at least three and at every positive scale.
This refutes the injectivity conjecture of Adamaszek, Frick and Vakili.
\end{abstract}

For a nonempty finite set $X\subset\mathbb R^n$ and a scale $r>0$, let
\[
 K=\VR(X;r)=\{\sigma\subseteq X:\operatorname{diam}(\sigma)<r\},
 \qquad
 S(X;r)=\bigcup_{\sigma\in K}\conv(\sigma),
\]
where distances are Euclidean. The vertex inclusion extends affinely on each
abstract simplex to a continuous map $p:|K|\to S(X;r)$, called the shadow map.
The complex is determined by its edges: a set of vertices spans a simplex
exactly when every pair is at distance less than $r$.

Chambers, de Silva, Erickson and Ghrist proved that the shadow map induces an
isomorphism on fundamental groups for planar samples
\cite[Theorem~3.1]{CDSEG}. The strict and non-strict scale conventions are
equivalent for these finite-sample statements by choosing a scale in a gap
between the finitely many distances \cite[Remark~1.4]{AFV}.
Adamaszek, Frick and Vakili proved surjectivity in three-dimensional space
\cite[Theorem~I]{AFV} and conjectured injectivity in all dimensions
\cite[Conjecture~7.2]{AFV}. The following example answers that fixed-sample
question negatively already in dimension three.

\begin{theorem}\label{thm:counterexample}
For every integer $n\ge3$ and every $r>0$, there is a six-point set
$X\subset\mathbb R^n$ such that $S(X;r)$ is contractible and
\[
 \pi_1(|\VR(X;r)|,v)\cong\mathbb Z
 \qquad\text{for every }v\in X.
\]
Consequently the affine shadow map is not injective on fundamental groups
at any sample base vertex.
\end{theorem}

\begin{proof}
We first take $n=3$ and $r=1$. Put
\[
 a=\frac9{20},\qquad b=\frac{501}{1000},\qquad s=\frac{309}{1000},
\]
and let $X=\{v_0,\ldots,v_5\}$, where
\begin{align*}
 v_0&=(a,s,0),&v_1&=(a,-s,0),&v_2&=(-b,0,0),\\
 v_3&=(b,0,0),&v_4&=(-a,0,s),&v_5&=(-a,0,-s).
\end{align*}
All six points are distinct. The following table lists every unordered pair;
the notation $ij$ denotes the pair $\{v_i,v_j\}$.
\begin{center}
\begin{tabular}{lll}
\toprule
Pairs & Squared distance & Comparison with $1$\\
\midrule
$01,45$ & $4s^2=95481/250000$ & $<1$\\
$02,12,34,35$ & $(a+b)^2+s^2=499941/500000$ & $<1$\\
$03,13,24,25$ & $(b-a)^2+s^2=49041/500000$ & $<1$\\
$04,05,14,15$ & $4a^2+2s^2=500481/500000$ & $>1$\\
$23$ & $4b^2=251001/250000$ & $>1$\\
\bottomrule
\end{tabular}
\end{center}
In particular, no distance equals the threshold. The only maximal simplices
of $K=\VR(X;1)$ are
\[
 [012],\quad[013],\quad[245],\quad[345].\tag{1}
\]
Indeed, there is no edge between $\{0,1\}$ and $\{4,5\}$, and the edge $23$
is missing. On each of the vertex sets $\{0,1,2,3\}$ and $\{2,3,4,5\}$,
this leaves precisely the two triangles displayed in (1).

Let $C\subset K$ be the four-cycle with ordered vertices $2,0,3,4$.
Define a map on abstract vertices by
\[
 f(1)=0,\qquad f(5)=4,\qquad f(i)=i\quad(i\in\{0,2,3,4\}).
\]
It sends the four triangles in (1) to the edges $[02],[03],[24],[34]$,
respectively, and extends to a simplicial retraction $f:|K|\to|C|$.
Moreover, for every simplex $\sigma$ in (1), the vertices of
$\sigma\cup f(\sigma)$ are still contained in $\sigma$. In the standard
barycentric realization, therefore,
\[
 H(u,t)=(1-t)u+t f(u),\qquad 0\le t\le1,
\]
stays in that same simplex. These affine formulas agree on shared faces,
and $H$ fixes $|C|$ pointwise. Thus $|K|$ strongly deformation retracts onto
the four-cycle. Its fundamental group is infinite cyclic
\cite[Theorem~1.7 and Proposition~1.17]{Hatcher}; a generator at $v_2$ is the
edge loop
\[
 \gamma=(v_2,v_0,v_3,v_4,v_2).\tag{2}
\]

The shadow has a different topology. Set
\[
 Q_1=\conv\{v_0,v_1,v_2\}\cup\conv\{v_0,v_1,v_3\},\qquad
 Q_2=\conv\{v_2,v_4,v_5\}\cup\conv\{v_3,v_4,v_5\}.
\]
By (1), the whole shadow is $Q_1\cup Q_2$. In the $xy$-plane the first
triangle has slices
$-b\le x\le a$, $|y|\le s(x+b)/(a+b)$, and the second has slices
$a\le x\le b$, $|y|\le s(b-x)/(b-a)$. Since $0<a<b$, their union is exactly
\[
 Q_1=\left\{(x,y,0):
 |y|\le\frac{s(x+b)}{a+b},\quad
 |y|\le\frac{s(b-x)}{b-a}\right\}.\tag{3}
\]
The two bounds agree at $x=a$; the first is smaller for $x\le a$ and the
second for $x\ge a$. Their nonnegativity forces $-b\le x\le b$.
The four linear inequalities in (3) show that $Q_1$ is a convex
quadrilateral. It contains the entire axis segment $[v_2,v_3]$.
The linear isometry $(x,y,z)\mapsto(-x,z,y)$ sends $Q_1$ to $Q_2$,
so $Q_2$ is also convex and contains that segment.

Both quadrilaterals contain the origin. Hence
\[
 h(w,t)=(1-t)w,\qquad w\in Q_1\cup Q_2,\quad 0\le t\le1,
\]
is a contraction of the shadow: a point stays in whichever convex
quadrilateral originally contains it. The projected loop $p(\gamma)$ is
therefore null-homotopic, whereas (2) is nontrivial in $|K|$.
The complex is connected, so its fundamental group is infinite cyclic at
every sample base vertex \cite[Proposition~1.5]{Hatcher}; the contractible
shadow has trivial fundamental group at all of them.

For $n\ge3$, append $n-3$ zero coordinates. For any prescribed $r>0$,
multiply the six points by $r$. Squared distances and the squared threshold
are both multiplied by $r^2$, giving exactly the same abstract complex,
while the shadow is carried to the new one by an isometric inclusion and
a positive scaling. This proves the theorem.
\end{proof}

The segment $[v_2,v_3]$ explains the lost loop. Its endpoints are too far
apart to be a Rips edge, but the segment lies in both planar parts of the
shadow. The two abstract triangulated disks meet only in the vertices
$v_2,v_3$; their shadows meet along the whole segment.

\paragraph{Exact verification.}
The accompanying standard-library Python code enumerates the exact rational
distances and all simplices, checks the carried vertex retraction, and checks
an integer simplicial cocycle taking value $1$ on (2). These finite checks
supplement the explicit homotopies above.

\begin{thebibliography}{9}
\bibitem{AFV}
Micha\l{} Adamaszek, Florian Frick and Adrien Vakili,
\emph{On homotopy types of Euclidean Rips complexes},
\href{https://arxiv.org/abs/1602.04131v1}{arXiv:1602.04131v1}, 2016.

\bibitem{CDSEG}
Erin W. Chambers, Vin de Silva, Jeff Erickson and Robert Ghrist,
\emph{Rips Complexes of Planar Point Sets},
\href{https://arxiv.org/abs/0712.0395v1}{arXiv:0712.0395v1}, 2007.

\bibitem{Hatcher}
Allen Hatcher, \emph{Algebraic Topology}, Cambridge University Press, 2002.
\href{https://math.uchicago.edu/~farb/algtop2023/Hatcher.pdf}{Original monograph PDF}.
\end{thebibliography}

\paragraph{AI-use statement.}
OpenAI Codex was used to develop and analyze the construction, check exact
arithmetic, and draft the manuscript and verification code. The mathematical
argument is given in full above. Independent mathematical review is requested.
\end{document}
